\documentclass[10pt,twocolumn]{article}
\usepackage[T1]{fontenc}
\usepackage[utf8]{inputenc}
\usepackage[margin=0.72in,top=0.85in,bottom=0.9in,columnsep=0.22in]{geometry}
\usepackage{mathptmx}           % Times Roman body + math
\usepackage{booktabs}
\usepackage{graphicx}
\usepackage{microtype}
\usepackage{amsmath}
\usepackage[colorlinks=true,linkcolor=black,citecolor=black,urlcolor=black]{hyperref}
\usepackage{titlesec}
\usepackage[font=small,labelfont=bf,labelsep=period,skip=4pt]{caption}

\graphicspath{{../results/figures/}}

%% ── Paragraph style ─────────────────────────────────────────────────────────
\setlength{\parindent}{1em}
\setlength{\parskip}{0pt}

%% ── Section headings (IEEE style) ──────────────────────────────────────────
\titleformat{\section}
  {\normalfont\normalsize\bfseries\centering}
  {\Roman{section}.}{0.5em}{\MakeUppercase}
\titleformat{\subsection}
  {\normalfont\normalsize\itshape}
  {\Alph{subsection}.}{0.5em}{}
\titlespacing*{\section}{0pt}{10pt plus 2pt minus 1pt}{4pt plus 1pt}
\titlespacing*{\subsection}{0pt}{7pt plus 1pt}{3pt plus 1pt}

%% ── Suppress default date ───────────────────────────────────────────────────
\date{}

\begin{document}

%% ── Title block (spans both columns) ───────────────────────────────────────
\twocolumn[{%
  \begin{center}
    {\large\bfseries
      Classical EEG Feature Robustness Under Subject-Aware Cross-Dataset Evaluation
    \par}
    \vskip 0.55em
    {\normalsize Barış Talar \quad \texttt{baristalar21@gmail.com}\par}
  \end{center}
  \vskip 0.5em
  \noindent\rule{\linewidth}{0.5pt}
  \vskip 0.45em
  {\small\noindent\textbf{Abstract}\textemdash
    Cross-dataset robustness remains a major barrier for motor-imagery EEG
    brain-computer interfaces. This study compares classical feature
    representations under zero-shot transfer between PhysioNet EEG Motor Imagery
    and BCI Competition IV Dataset 2a. The revised analysis corrects key protocol
    weaknesses in the initial benchmark: time-domain features are strengthened,
    subject identifiers are preserved, within-dataset estimates use subject-grouped
    cross-validation, feature dimensions are controlled by PCA, and cross-dataset
    uncertainty is bootstrapped over target subjects. Under this stricter protocol,
    no classical representation provides reliable transfer. PCA-matched
    cross-dataset balanced accuracy averaged across models was 0.681, 0.645, and
    0.604 for Time Domain, Band Power, and FFT in the PhysioNet-to-BCI2a direction,
    but only 0.559, 0.545, and 0.557 in the reverse direction. Pairwise
    target-subject tests show small effects: Time Domain exceeds Band Power by
    0.014 ($p=0.0006$), while Time Domain versus FFT is not significant
    ($p=0.179$). The previous FFT-over-Band-Power-over-Time-Domain claim is not
    supported.
  \par}
  \vskip 0.45em
  \noindent\rule{\linewidth}{0.4pt}
  \vskip 1.1em
}]

\section{Introduction}
Motor-imagery EEG classification is often evaluated within a single dataset,
where subjects, acquisition hardware, and protocol details are partly shared
between train and test data. Such evaluations can overstate deployment
readiness. A useful BCI model must tolerate changes in subjects, electrodes,
amplifiers, and task execution.

The initial version of this benchmark asked whether frequency-domain features
were more robust than time-domain statistics. That framing was too strong:
pooled trial-wise validation mixed subjects across folds, the time-domain
baseline used only mean and variance after bandpass filtering, feature
dimensions were imbalanced, and the headline ranking relied on very small
normalised-gap differences. This revision treats the work as a controlled
negative benchmark: it tests whether any classical representation remains
reliable under subject-aware cross-dataset evaluation.

\section{Methods}
Two MOABB datasets were used. PhysioNet EEG Motor Imagery contributes 109
subjects and 4,918 retained left/right imagery trials after preprocessing.
BCI Competition IV Dataset 2a contributes 9 subjects and 2,592 retained
left/right trials. Recordings were reduced to the 22 BCI2a channels, resampled
to 250 Hz, filtered from 4--40 Hz, common-average referenced, and epoched from
0.5--2.5 s after cue onset. Trial-level subject, session, run, event, and label
metadata are written to CSV files.

Three representations were compared. Time Domain contains 154 descriptors:
mean, standard deviation, RMS, line length, zero-crossing rate, Hjorth mobility,
and Hjorth complexity for each channel. Band Power contains log Welch power in
theta, mu/alpha, beta, and low-gamma bands for each channel. FFT contains log
magnitude spectra from 4--40 Hz interpolated to 36 bins per channel. The
primary comparison uses PCA inside the model pipeline to match all feature
sets to 88 dimensions.

Models were Logistic Regression, linear SVM, and Random Forest. The primary
within-dataset protocol uses stratified subject-grouped folds. Trial-wise
stratified CV is retained only as a diagnostic legacy baseline. Cross-dataset
evaluation trains on all source subjects and tests zero-shot on all target
subjects in both directions. Cross-dataset confidence intervals are obtained by
bootstrapping target subjects. Feature comparisons use paired sign-flip
permutation tests over target-subject balanced accuracies.

\section{Results}
Table~\ref{tab:within} reports PCA-matched subject-grouped within-dataset
balanced accuracy. The strengthened time-domain representation is no longer
degenerate and is strongest within both datasets.

\begin{table}[t]
\centering
\caption{Subject-grouped within-dataset balanced accuracy.}
\label{tab:within}
\small
\begin{tabular}{@{}llcc@{}}
\toprule
Feature & Model & PhysioNet & BCI2a \\
\midrule
Time & LogReg & 0.613 & 0.676 \\
Time & SVM    & 0.616 & 0.676 \\
Time & RF     & 0.579 & 0.645 \\
\midrule
Band & LogReg & 0.603 & 0.637 \\
Band & SVM    & 0.603 & 0.634 \\
Band & RF     & 0.563 & 0.588 \\
\midrule
FFT  & LogReg & 0.569 & 0.592 \\
FFT  & SVM    & 0.570 & 0.592 \\
FFT  & RF     & 0.524 & 0.559 \\
\bottomrule
\end{tabular}
\end{table}

Table~\ref{tab:cross} reports zero-shot cross-dataset balanced accuracy with
target-subject bootstrap intervals. Transfer is asymmetric: training on the
larger PhysioNet source and testing on BCI2a is easier than training on 9
BCI2a subjects and testing on 109 PhysioNet subjects.

\begin{table}[t]
\centering
\caption{Cross-dataset balanced accuracy with 95\% target-subject bootstrap
  intervals.}
\label{tab:cross}
\footnotesize
\setlength{\tabcolsep}{4pt}
\begin{tabular}{@{}llcc@{}}
\toprule
Feature & Model & Phys $\to$ BCI & BCI $\to$ Phys \\
\midrule
Time & LogReg & 0.699 {[}0.639, 0.774{]} & 0.579 {[}0.559, 0.599{]} \\
Time & SVM    & 0.700 {[}0.641, 0.774{]} & 0.577 {[}0.557, 0.598{]} \\
Time & RF     & 0.644 {[}0.575, 0.722{]} & 0.520 {[}0.506, 0.534{]} \\
\midrule
Band & LogReg & 0.659 {[}0.580, 0.750{]} & 0.558 {[}0.541, 0.576{]} \\
Band & SVM    & 0.654 {[}0.573, 0.748{]} & 0.554 {[}0.537, 0.571{]} \\
Band & RF     & 0.623 {[}0.558, 0.699{]} & 0.525 {[}0.512, 0.537{]} \\
\midrule
FFT  & LogReg & 0.626 {[}0.573, 0.678{]} & 0.573 {[}0.556, 0.590{]} \\
FFT  & SVM    & 0.623 {[}0.570, 0.676{]} & 0.572 {[}0.555, 0.589{]} \\
FFT  & RF     & 0.563 {[}0.525, 0.608{]} & 0.526 {[}0.512, 0.540{]} \\
\bottomrule
\end{tabular}
\end{table}

The target-subject paired permutation tests give: Time Domain minus Band Power
= 0.014 ($p=0.0006$), Time Domain minus FFT = 0.0066 ($p=0.179$), and Band
Power minus FFT = -0.0075 ($p=0.146$). These are small effects, not a strong
feature hierarchy.

\section{Discussion}
The revised benchmark changes the interpretation. The old near-chance
PhysioNet result was partly caused by a weak time-domain baseline and pooled
analysis choices. With common-average referencing and richer time-domain
descriptors, PhysioNet subject-grouped accuracy rises to about 0.61 for linear
models, and many individual subjects exceed 0.70 in within-subject evaluation.

The main conclusion remains conservative. Classical features can recover
motor-imagery signal in subject-specific or subject-aware settings, but they do
not provide robust zero-shot transfer across datasets. The normalised
generalisation gap is also fragile: in the PhysioNet-to-BCI2a direction, cross
accuracy can exceed the source grouped baseline because the target dataset is
smaller and has stronger responders. Direct cross-dataset accuracy,
target-subject intervals, and paired tests are therefore more informative than
a single averaged gap.

\section{Limitations}
Only two motor-imagery datasets are studied. The PCA-matched comparison
controls dimensionality but cannot make representations semantically
equivalent. The model set is deliberately classical and does not yet include
CSP, FBCSP, Riemannian classifiers, deep learning, or domain adaptation. These
baselines should be added before making broad claims about EEG transfer
learning.

\section{Conclusion}
After subject-aware evaluation, fairer time-domain features, dimensionality
matching, and target-subject uncertainty estimates, the original FFT robustness
claim is not supported. No tested classical feature representation provides
reliable zero-shot cross-dataset motor-imagery EEG transfer. Future work should
add CSP/FBCSP, Riemannian geometry, and domain-adaptation baselines before
positioning this as a mature BCI transfer benchmark.

\begin{thebibliography}{9}

\bibitem{brunner2008}
C.~Brunner, R.~Leeb, G.~Mueller-Putz, A.~Schloegl, and G.~Pfurtscheller,
``BCI Competition 2008 -- Graz data set A,'' 2008.

\bibitem{goldberger2000}
A.~L. Goldberger et al.,
``PhysioBank, PhysioToolkit, and PhysioNet,''
\emph{Circulation}, vol.~101, no.~23, pp.~e215--e220, 2000.

\bibitem{jayaram2018}
V.~Jayaram and A.~Barachant,
``MOABB: Trustworthy algorithm benchmarking for BCIs,''
\emph{Journal of Neural Engineering}, vol.~15, no.~6, 2018.

\bibitem{jayaram2016}
V.~Jayaram et al.,
``Transfer learning in brain-computer interfaces,''
\emph{IEEE Computational Intelligence Magazine},
vol.~11, no.~1, pp.~20--31, 2016.

\bibitem{lotte2018}
F.~Lotte et al.,
``A review of classification algorithms for EEG-based brain-computer
interfaces: A 10-year update,''
\emph{Journal of Neural Engineering}, vol.~15, no.~3, 2018.

\bibitem{pfurtscheller2001}
G.~Pfurtscheller and C.~Neuper,
``Motor imagery and direct brain-computer communication,''
\emph{Proceedings of the IEEE},
vol.~89, no.~7, pp.~1123--1134, 2001.

\bibitem{schalk2004}
G.~Schalk et al.,
``BCI2000: A general-purpose brain-computer interface system,''
\emph{IEEE Transactions on Biomedical Engineering},
vol.~51, no.~6, pp.~1034--1043, 2004.

\end{thebibliography}

\end{document}
